\documentclass[10pt,a4paper]{amsart}
\usepackage[margin=19mm]{geometry}
\usepackage{amsmath,amssymb,mathtools}
\usepackage[scr=rsfs,cal=euler]{mathalfa}
\usepackage{xfrac}
\usepackage[unicode,hidelinks,bookmarks=false]{hyperref}
\newcommand{\ie}{i.e.\ }
\newcommand{\cf}{cf.\ }
\newcommand{\resp}{respectively\ }
\newcommand{\Subsection}{\subsection}
\providecommand{\nobreakdash}{\nobreak\hbox{-}}
\setlength{\emergencystretch}{2em}
\multlinegap0pt
\usepackage{amssymb}
\usepackage{xcolor}
\usepackage{tikz}
\newtheorem{theorem}{Theorem}
\def\r{\mathbb R}
\DeclareMathOperator{\sech}{sech}
\def\v{\mathcal V}
\title{Prescribed Angle Surfaces Associated with Torse-Forming Vector Fields in Riemannian Manifolds}
\author{Muhittin Evren Ayd{\i}n, Esra Dilmen, B\"u\c{s}ra Karakaya}
\date{Source: arXiv:2605.09553v1 (CC BY 4.0)}
\begin{document}
\maketitle

\noindent\emph{Source and licence.} Prescribed Angle Surfaces Associated with Torse-Forming Vector Fields in Riemannian Manifolds. Version arXiv:2605.09553v1 (CC BY 4.0), distributed by arXiv under the Creative Commons Attribution 4.0 International licence, http://creativecommons.org/licenses/by/4.0/, as recorded in the arXiv licence metadata for this exact version and shown on the abstract page https://arxiv.org/abs/2605.09553v1. Full licence notice: \texttt{licenses/arxiv-2605.09553-CC-BY-4.0.txt}.\par\smallskip

\noindent\textit{Editorial scope.} Counted unit: the article's theorem classtotgeo (source0.tex lines 470--490). The article's complete original proof of it is delivered with it (source0.tex lines 491--534, one \texttt{proof} environment of 44 lines). No mathematics is written, repaired or re-derived: the statement and the proof are the article's own lines, in the article's own notation, and no retained line is substituted or re-flowed.  Retained uncounted as the source-local context the proof needs: lines 428--430.  Nothing was omitted from any retained span, so the package carries no omission record.  No retained line contains a citation, so the wrapper carries no bibliography and no bibliography entry.  No retained line contains a figure environment, an image-inclusion command or drawing code, so no image is delivered and the package declares no asset dependency.  Editorial formatting only, in addition: an AMS \texttt{amsart} preamble that restates the article's own declarations and macros used by the retained lines, namely the environments \texttt{theorem} on separate counters, together with the standard packages those lines need.  The retained environments print in this document as Theorem 1 in order of appearance, whereas the article prints the same items with the numbering of its own full document.  The complete native source of the article as distributed in the arXiv e-print for this exact version is bundled unchanged as \texttt{sources/200233/source0.tex}.
\begin{equation}\label{tg1}
\theta(u)=\cos^{-1}\big( \sech F(u) \big), \quad K_{\mathrm{int}}=-\big(F''(u)\coth F(u)+F'(u)^2\big).
\end{equation}

\begin{theorem}\label{classtotgeo}
Let $\Sigma$ be a totally geodesic PA surface with $(\v,\theta)$, where $\v$ is an anti-torqued vector field with nonzero potential function $f$ and $\theta$ is non-constant. The intrinsic curvature $K_{\mathrm{int}}$ is some constant $K_0$ if and only if there is a coordinate system $(u,v)$ on $\Sigma$ such that 
\begin{equation}\label{ODEtg0}
F(u)=
\begin{cases}
\displaystyle
\sinh^{-1}\!\left(
\sqrt{\frac{c_1-K_0}{-K_0}}\,
\sinh\big(\pm\sqrt{-K_0}(u+c_2)\big)
\right) , & K_0<0, \, c_1\geq 0,\\[14pt]
\displaystyle
\sinh^{-1}\big(\pm\sqrt{c_1}\,u+c_2\big), & K_0=0, \, c_1>0,\\[10pt]
\displaystyle
\sinh^{-1}\!\left(
\sqrt{\frac{c_1-K_0}{K_0}}\,
\sin\big(\sqrt{K_0}(u+c_2)\big)
\right), & 0<K_0<c_1, 
\end{cases}
\end{equation}
where $c_1,c_2\in \r$  and the potential function is given by $f=F'$.
\end{theorem}

\begin{proof}
Assume that $K_{\mathrm{int}}=K_0$. Then, by \eqref{tg1}, we obtain
\begin{equation}\label{ODEtg}
F''\coth F+F'^2+K_0=0.
\end{equation}
Since $f\neq 0$, the function $F$ is non-constant. If $F''=0$, then \eqref{ODEtg} reduces to $F'^2+K_0=0$, which implies that $K_0<0$ and $F'(u)=\pm\sqrt{-K_0}$. This corresponds to a particular solution of the form $F(u)=\pm\sqrt{-K_0}u+c$ in \eqref{ODEtg0}, where $c_1=0$ and $c_2=c$.

We now assume that $F''\neq 0$ and proceed to solve \eqref{ODEtg}. Since $F''\neq 0$, we may treat the function $F'$ as a variable. Therefore, we set $F'=P$. By the chain rule, differentiating this equality yields
$$F''=\frac{dF}{du}\frac{dP}{dF}=P\frac{dP}{dF}.$$ 
Substituting into \eqref{ODEtg}, we conclude
$$
P\frac{dP}{dF}\coth F+P^2+K_0=0.
$$
We change the variable again and set $Q=P^2$. Differentiating, we obtain $\frac{dQ}{dF}=2P\frac{dP}{dF}$, and thus
$$
\frac{1}{Q+K_0}\frac{dQ}{dF}=-2\tanh F.
$$
Integrating,

\begin{equation}\label{ODEtg001}
P^2=c_1\sech^2F-K_0,
\end{equation}
where $c_1>0$ is some constant.

Finally, we use another change of variable. Introduce $S(u)=\sinh F(u)$. Then $S'=F'\cosh F$. Multiplying \eqref{ODEtg001} by $\cosh^2 F=1+\sinh^2 F$, we obtain
\begin{equation}\label{ODEtg1}
S'^2=c_1-K_0(1+S^2)=c_1-K_0-K_0S^2.
\end{equation}
The solutions depend on the sign of $K_0$.

\textbf{Case $K_0<0$.} Solving \eqref{ODEtg1}, we obtain
$$
S(u)=\sqrt{\frac{K_0-c_1}{K_0}}\sinh \big (\pm\sqrt{-K_0}u+c_2 \big ), \quad c_2\in \r.
$$
Since $S(u)=\sinh F(u)$, this yields the expression for $F(u)$ in the case $K_0<0$.

\textbf{Case $K_0=0$.} Equation \eqref{ODEtg1} reduces to $S'^2=c_1$ and so $S(u)=\pm \sqrt{c_1}u+c_2$, $c_2\in \r.$ Thus $F(u)=\sinh^{-1}(\sqrt{c_1}u+c_2)$. 

\textbf{Case $K_0>0$.} In this case, \eqref{ODEtg1} implies $c_1>K_0$ and we derive
$$
S(u)=\sqrt{\frac{c_1-K_0}{K_0}}\sin \big (\sqrt{K_0}u+c_2 \big ), \quad c_2\in \r.
$$
Again, using $S(u)=\sinh F(u)$, we obtain the required expression for $F(u)$ in the case $K_0>0$.
\end{proof}

\end{document}
